\documentclass[../main.tex]{subfiles}
\begin{document}
%==========================================TIÊU ĐỀ TẬP========================================
\begin{table}[h]
  \begin{adjustwidth}{-1cm}{}
    \begin{tabular}{>{\centering\arraybackslash}p{6cm}|>{\raggedright\arraybackslash}p{12cm}}
      \multirow{5}{6cm}
      % Cột bên trái: ÉPISODE và số tập
      {\\[-40pt]
        \centering
        {\fontsize{20pt}{20pt}\selectfont{\outlinetext[1]{eptype}{white}{\flaregothic ÉPISODE}}\\[12pt]
          {\fontsize{72pt}{72pt}\selectfont{\outlinetext[1]{epnum}{white}{\burbank 169}}}\\[6pt]       % use for normal/special episodes
          % {\fontsize{36pt}{36pt}\selectfont{\outlinetext[1]{epnum}{white}{\burbank SÁU}}}\\[6pt]     % use for custom episodes
          {\fontsize{20pt}{20pt}\selectfont{\outlinetext[1]{epnum}{white}{\cafeteria (13-10)}}}\color{black}}
        \\[18pt]
        % {\fontsize{14pt}{14pt}\selectfont{\outlinetext[1]{epattr}{white}{\flaregothic [I]}}}\\
        % {\fontsize{18pt}{18pt}\selectfont{\outlinetext[1]{epattr}{white}{\flaregothic [II]}}}\\
        \color{black}
      }
       &
      % Dòng thứ nhất: tên tập tiếng Nhật
      {\fotskip\fontsize{18pt}{18pt}\selectfont セミうっせぇ！マジで！}                                                                                             \\[6pt]
      % Dòng thứ hai: tên tập tiếng Anh
       & {\cafeteria\fontsize{18pt}{18pt}\selectfont When The Cicadas Fly}                                                                               \\[4pt]
      % Dòng thứ ba: tên tập tiếng Việt
       & {\vnmsans\fontsize{18pt}{18pt}\selectfont Dàn đồng ca mùa hạ}                                                                                   \\[8pt]
      % Dòng thứ tư: Nhân vật xuất hiện
       & {\caviar{\fontsize{14pt}{14pt}\selectfont Nhân vật: \emp{kuon2} \emp{kaiho2} \emp{suzuno2} \emp{fubuki} \emp{haato} \emp{pekora} \emp{koyori}}}
      \\[8pt]
      % Dòng cuối cùng: Thời gian diễn ra của tập (thường sẽ là ngày phát hành)
       & {\sen\fontsize{16pt}{16pt}\selectfont Mốc thời gian: 07/08/2022}                                                                                \\[18pt]
    \end{tabular}\\[8pt]
  \end{adjustwidth}
\end{table}
%=========================================TRUYỆN CHÍNH========================================
\color{black}

\txtbx[2]{
  {\color{vol} \uline{Tóm tắt:}} Tiếng ve sầu mùa hạ râm ran nhức óc đẩy hai cô nàng tai thú Fubuki và Pekora vào cảnh khổ sở dở khóc dở cười, trong khi Koyori lại bày trò tự uống thuốc ``làm thủng màng nhĩ'' để rồi khốn đốn khi thuốc giải có tác dụng. Cơn ác mộng thực sự bùng nổ khi Akai Haato thả nguyên một ổ ve sầu nuôi trong rừng bay tán loạn khắp văn phòng. Giữa lúc Kuon, Kaiho và cô em gái Suzuno vất vả đuổi bắt đàn côn trùng thì Fubuki, Pekora và Koyori lại tự giam mình trong phòng thu âm không điều hòa để rồi bị mức nhiệt gay gắt giữa hè thiêu đốt đến tan chảy.
}

\pov{Hikawa Kuon}

Khi nhắc tới mùa hè, ta thường liên tưởng đến điều gì?

Bầu trời xanh ngắt không một gợn mây, ánh nắng rực rỡ chói chang, hay vị ngọt lịm mát lạnh của những que kem trái cây xua tan đi cái nóng nực oi ả?

Nếu nói về thứ thanh âm đặc trưng tạo nên linh hồn của một mùa hè xứ Phù Tang đúng nghĩa, thì chắc chắn đó phải là tiếng ve sầu. Thứ côn trùng nhỏ bé ấy từng là nguồn cảm hứng bất tận cho biết bao áng thơ ca, những khúc nhạc trong trẻo của thời cắp sách tới trường. Chính vì thế, người ta vẫn thường ví von tiếng ve kêu râm ran như một bản giao hưởng êm đềm, một giai điệu đầy hoài niệm báo hiệu những ngày hè rực rỡ đã sang.

Thế nhưng, ở chốn văn phòng \hll\ này, thứ thanh âm ấy dường như lại\dots\ chẳng hề được hoan nghênh cho lắm. Thậm chí nếu nói thẳng ra, nó còn là một cơn ác mộng kinh hoàng.

Cũng phải thôi, chúng quá ồn ào. Và đối với các cô gái ở Thành phố Điện tử này, sự ồn ào ấy lại càng trở thành nỗi khiếp đảm tột độ. Nhất là khi có ai đó cả gan đem nguyên một ổ ve sầu thả rông ngay giữa trụ sở\dots

\begin{center}
  \uline{OE OE\dots\ OE OE\dots\ OE OE\dots}
\end{center}

Lại là một ngày ngột ngạt và oi bức đến phát điên.

Cái bức bối ấy không chỉ bắt nguồn từ đợt nắng nóng đỉnh điểm của những ngày đầu tháng Tám\dots\ mà còn bởi những tiếng ve sầu râm ran đinh tai nhức óc cứ dội thẳng vào màng nhĩ suốt từ sáng tinh mơ tới giờ. Chúng ồn ào đến mức khiến bất kỳ ai nghe thấy cũng phải phát cáu.

\begin{center}
  \uline{OE OE\dots\ OE OE\dots\ OE OE\dots}
\end{center}

Fubuki và Pekora lúc này đang trong tình cảnh vô cùng khốn khổ. Cả hai cô nàng tai thú vừa dùng hai tay bịt chặt lấy đôi tai nhạy cảm của mình trong sự tuyệt vọng, vừa kiễng gối cúi rạp người xuống gần sát mặt sàn, nghiến răng kèn kẹt chịu đựng những đợt sóng âm như khoan thủng não bộ.

\img{13-10a}

Fubuki và Pekora đồng thanh hét toáng lên: ``\textbf{Ồn quá đi mất!! Tai của tôi sắp nổ tung rồi!!}''

Kaiho ngồi bên cạnh tôi trên chiếc ghế sô-pha dài, tay cầm cốc trà lúa mạch mát lạnh, chép miệng thở dài: ``\dots Trông tội nghiệp hai đứa nhỏ ghê chưa.''

Thực ra mà nói, việc cô em gái song sinh kiêm Tổng quản lý chi nhánh nam \hls\ xuất hiện ở văn phòng \hll\ của tôi hôm nay hoàn toàn có lý do chính đáng. Hôm nay là ngày Chủ nhật, bên chi nhánh nam không có lịch ghi hình hay sự kiện đột xuất nào, các anh tài lớn đều tự luyện tập cá nhân hoặc rủ nhau ra ngoài dạo phố.

Nhưng quan trọng hơn cả là chuyện gia đình. Suốt ba tuần lễ nay, kể từ sau biến cố đau thương vào sáng sớm ngày mười sáu tháng Bảy - ngày mà ba mẹ con nhà Chó (Miyako, Toa, Kiri) bỗng dưng bốc hơi không một dấu vết sau vụ trộm mộ kỳ quái ở xưởng cơ khí ngoại ô - bầu không khí trong khu ký túc xá phố Anwa lúc nào cũng chùng xuống, nặng trĩu một nỗi u hoài. Người chịu tổn thương và buồn bã nhất không ai khác chính là Suzuno, cô em gái thứ tư vốn dĩ rất mực yêu thương và coi ba chú cún như những người bạn tri kỷ gắn bó. Thấy con bé cứ ngày ngày ủ rũ, hết giam mình trong phòng đọc sách lại thơ thẩn ra ban công nhìn về phía chân trời xa xăm, Kaiho không đành lòng để em gái chìm mãi trong u uất. Thế là sáng nay, nhân một ngày rảnh rỗi, cô nàng liền chủ động dắt díu Suzuno ra ngoài đổi gió, tiện tay mang theo một âu chè đậu đỏ ướp lạnh mẹ nấu sang tiếp tế cho tôi, sẵn tiện ghé thăm xem ông anh trai song sinh đang xoay xở công việc ở trụ sở bên này ra sao.

Nào ngờ đâu, vừa đặt chân tới văn phòng chưa được bao lâu, cả ba anh em chúng tôi đã đụng ngay đợt cắt điện bảo trì lưới điện luân phiên của khu phố, lại còn phải chịu trận trước cơn bão âm thanh kinh hoàng từ dàn đồng ca ve sầu này.

Tôi nhấp một ngụm nước trà, nhún vai: ``Biết làm sao được? Thính giác của loài cáo với loài thỏ nhạy bén hơn tai người thường gấp nhiều lần, hai em ấy lãnh đủ là điều tất yếu thôi\dots''

``Chưa kể chúng ta đã quá quen với mấy thứ tiếng ồn còn khủng khiếp hơn thế này nhiều lần rồi\dots'' cô em song sinh của tôi chống cằm lẩm bẩm.

Tôi liếc xéo sang: ``Không cần phải đem cái tiền sử tai ương chấn động màng nhĩ ấy ra để khoe mẽ đâu.''

Tôi và Kaiho vốn đã trải qua đủ thứ âm thanh bão táp trên cõi đời này, nên mấy tiếng ve sầu này đối với chúng tôi quả thực chẳng bõ bèn gì.

Bên cạnh tôi, Suzuno đang ngồi lặng lẽ tựa đầu vào thành ghế. Đôi mắt dị sắc tĩnh lặng như mặt hồ không một gợn sóng, khẽ dõi theo hai bậc tiền bối đang quằn quại dưới sàn. Dù đã được chị gái kéo ra khỏi bốn bức tường ngột ngạt ở ký túc xá, nhưng nỗi nhớ thương ba người bạn nhỏ bốn chân dường như vẫn luôn đè nặng trong tâm trí con bé. Tôi khẽ đưa tay xoa nhẹ lên mái tóc mềm mại của em gái, truyền cho con bé chút hơi ấm an ủi giữa những ngày lòng người còn trĩu nặng.

Đúng lúc đó, một giọng nói hồ hởi, lanh lảnh quen thuộc bất ngờ cất lên: ``Em có một loại thuốc mà có thể làm thủng màng nhĩ của mọi người ngay tức khắc luôn đấy ạ!''

``Koyo-chan!?'' / ``Koyori-chan!?''

Cả tôi và Kaiho đồng loạt ngoảnh lại. Đập vào mắt chúng tôi là nàng Sói đồng cỏ, trên tay đang hớn hở cầm một chiếc lọ thủy tinh chứa thứ dung dịch màu xanh nước biển lấp lánh đầy ma quái.

\img{13-10b}

Nhưng nghe xong cái công dụng rợn người mà em ấy vừa dõng dạc tuyên bố, hai khóe miệng tôi giật giật không ngừng. Thuốc làm thủng màng nhĩ á!? Lại còn khoe khoang một cách tự hào như thế nữa chứ!

Koyori cười toe toét, khua khua chiếc lọ: ``Anh chị Tổng quản lý có vẻ như không bị tác động nên chắc không cần dùng ha\~{}''

Kaiho tròn mắt, khóe môi khẽ nhếch lên một nụ cười gượng gạo: ``Lọ thuốc đó trông lạ mắt đấy, cơ mà\dots\ Uống cái thứ đó vào thì chẳng phải sẽ điếc đặc vĩnh viễn luôn rồi hay sao?''

Tôi cũng gật đầu phụ họa: ``Với lại em cất công nghiên cứu ra cái thứ quái thai ấy để làm chi vậy?''

Koyori chớp chớp mắt, vênh mặt tự hào: ``Em đã biết trước là sẽ có ngày nó phát huy tác dụng mà! Hai chị có muốn dùng thử hông ạ?''

Dù tiếng ve sầu gầm rú bên ngoài át đi phần lớn âm thanh thông thường, nhưng với năng lực thần giao cách cảm và thính lực đặc thù của những kẻ dị biệt, chúng tôi vẫn có thể nghe rõ từng lời em ấy nói.

Thế nhưng, vừa dứt lời mời chào, ánh mắt của cô nàng Sói đồng cỏ vô tình chạm phải Suzuno đang ngồi bên cạnh tôi.

Đôi mắt hai màu của Suzuno khẽ ngước lên, bình thản nhìn chằm chằm vào chiếc lọ trên tay nàng Sói.

Ký ức kinh hoàng từ buổi chiều định mệnh ở ngày em ấy ra mắt lập tức dội về trong tâm trí Koyori như một cơn lốc: hình ảnh cô bé cấp hai nhà Hikawa điềm đạm liệt kê hàng loạt điều luật an toàn phòng thí nghiệm, đe dọa tước giấy phép nghiên cứu và nắn nót ghi lại từng hành vi bất thường vào cuốn sổ tay bệnh án!

Cô nàng giật bắn mình như vừa chạm tay vào dây điện cao thế, cả người cứng đờ, hai tai sói hồng lập tức cụp rạp xuống sát đầu. Em ấy vội vàng giấu biến chiếc lọ thuốc ra sau lưng, mồ hôi hột toát ra đầm đìa, miệng lắp bắp thanh minh: ``S-Suzuno-chan\dots! Em đừng hiểu lầm nhé! C-Cái này hoàn toàn là nghiên cứu thực nghiệm phi thương mại thôi, tuyệt đối không có mục đích gây hại cho cộng đồng đâu ạ\dots!''

Suzuno chỉ khẽ chớp mắt, cất giọng đều đặn, lễ phép: ``Theo quy chuẩn an toàn sinh học thì việc tự ý phá hủy cơ quan thính giác bằng hóa chất không qua kiểm định là hành vi bị nghiêm cấm đấy ạ, Koyori-san.''

``H-Hic\dots! Chị biết rồi mà\dots!'' Koyori co rúm người lại, mặt mày tái mét.

Trong khi đó, Fubuki và Pekora dù đang đau đớn bịt tai cũng không thể nhịn nổi sự phẫn nộ:

Nàng Thỏ gắt lên: ``Đừng có mà tiếp thị mấy thứ tự hủy hoại bản thân ấy nữa, peko!''

Nàng Cáo trắng cũng nghiến răng: ``Koyorin, bộ tiếng ồn này không làm phiền em chút nào sao\dots!?''

Đáp lại sự bức xúc của hai bậc tiền bối chỉ là một cái nghiêng đầu đầy ngơ ngác từ nàng Sói. Một dấu hỏi chấm to tướng như muốn nhảy múa trên đỉnh đầu em ấy.

Pekora giậm chân bình bịch xuống sàn, gào to hơn: ``\textbf{Em có nghe thấy chị nói gì không hả?! Đừng có đem ba cái loại thuốc vớ vẩn ấy ra quảng cáo nữa đi!}''

Fubuki hét lớn: ``\textbf{Em thực sự không thấy nhức óc vì tiếng ve sầu này ư, Koyorin\dots!?}''

Kaiho cũng cất giọng sang sảng: ``\textbf{Sói đồng cỏ ơi\~{}! Có nghe thấy mọi người đang nói gì không đấy\~{}?}''

Dù ba người họ có gân cổ lên gào thét đến lạc cả giọng, Pekora còn đập tay xuống bàn bôm bốp để thu hút sự chú ý, nhưng Koyori vẫn cứ trơ mắt ếch nhìn lại. Trên đầu em ấy lúc này không còn là dấu hỏi nữa, mà biến thành một chuỗi ba dấu chấm lơ lửng đúng nghĩa đen.

Tôi nhíu mày. Thính lực của loài sói đồng cỏ thậm chí còn nhạy bén hơn cả người thường, cớ sao em ấy lại không nghe thấy gì cơ chứ?

Khoan đã\dots\ ``Thuốc làm thủng màng nhĩ''?

Tôi và Pekora đồng loạt đưa mắt nhìn sang chiếc lọ dung dịch màu xanh lục mà Koyori vừa đặt hờ hững trên bàn lúc nãy. Mực nước bên trong đã vơi đi một nửa.

\img{13-10c}

Đúng như dự đoán tồi tệ nhất của chúng tôi.

Pekora hoảng hốt: ``Không lẽ nào\dots\ em ấy đã dùng nó rồi sao, peko!?''

Tôi thở hắt ra một hơi, day day sống mũi: ``Em ấy tự nốc bố nó rồi chứ còn gì nữa\dots''

Kaiho tặc lưỡi: ``Chả trách sao nãy giờ như người điếc đặc.''

Suzuno khẽ lắc đầu, lẩm bẩm ghi nhận: ``Tự lấy bản thân ra làm vật thí nghiệm phá hủy màng nhĩ mà không chuẩn bị trước thiết bị bảo hộ\dots\ Thật là thiếu cẩn trọng quá.''

Fubuki bất lực buông thõng hai tay: ``Ta cứ đứng đây gào thét thêm nữa thì em ấy cũng có hiểu chúng ta nói gì đâu\dots''

Dĩ nhiên, mọi lời chúng tôi nói đều hoàn toàn trôi tuột qua tai Koyori. Nhìn thấy khẩu hình miệng của mọi người chuyển động liên hồi mà bản thân chẳng tiếp nhận được mảy may thanh âm nào, nàng Sói hồng khoanh tay trầm ngâm suy nghĩ, tự hỏi vì sao mình lại chẳng nghe thấy gì.

Sau một thoáng vận dụng hết công suất chất xám ít ỏi, em ấy búng tay cái tách, vẻ mặt bừng sáng như vừa phát minh ra chân lý mới: ``À đúng rồi ha! Cái màng nhĩ của mình bị chọc thủng mất rồi còn đâu!''

Cặp tai sói màu hồng vểnh lên đầy phấn khởi. Tôi thở dài ngao ngán: ``Tốc độ truyền tín hiệu của não bộ em ấy có vấn đề nghiêm trọng rồi đấy\dots'' còn Kaiho nhăn nhó: ``Thì chúng ta cũng mất một lúc mới nhận ra còn gì.''

Koyori hí hửng lục lọi túi áo, lôi ra một chiếc lọ thủy tinh khác chứa dung dịch màu xanh lam trong suốt: ``Phải phục hồi lại nó ngay thôi!''

Nói đoạn, em ấy ngửa cổ tu sạch một hơi không chừa lại giọt nào. Trời ạ, nếu đã có thuốc giải điều chế sẵn thì sao không uống ngay từ đầu cho thiên hạ nhờ!

Lọ thuốc chữa màng nhĩ quả nhiên phát huy tác dụng với tốc độ thần kỳ. Nhưng đó cũng chính là lúc bi kịch thực sự bắt đầu.

Ngay khoảnh khắc thính giác được phục hồi 100\%, toàn bộ cường độ âm thanh của dàn đồng ca ve sầu bên ngoài ập thẳng vào hai tai sói của Koyori như một quả bom nổ chậm.

\begin{center}
  \uline{OE OE\dots\ OE OE\dots\ OE OE\dots!!}
\end{center}

Sắc mặt nàng Sói lập tức chuyển từ hồng hào sang trắng bệch, rồi xanh lét như tàu lá chuối. Em ấy ôm chặt lấy đầu, thét lên xé lòng trong đau đớn tột cùng: ``\textbf{Á á á á! Tai của tôi!! Cứu tôi với!!}''

\img{13-10d}

Pekora chống nạnh: ``Đấy, biết ngay mà! Tự làm tự chịu nhé, peko!''

Kaiho thấy vậy lại tỏ vẻ tiếc nuối: ``Tự dưng giờ chị lại muốn uống thử cái lọ làm thủng màng nhĩ ban nãy ghê\dots''

Fubuki trợn tròn mắt: ``Ủa, em tưởng chị và Kuon-san chịu được tiếng ồn cơ mà!?''

Kaiho nhún vai: ``Thì chịu được, nhưng nghe mãi cũng phiền. Với lại em ấy có sẵn thuốc giải kìa, lo gì?''

Suzuno khẽ kéo nhẹ vạt áo Kaiho, nghiêm túc khuyên nhủ: ``Onee-sama, đừng dại dột làm theo ạ. Thuốc phục hồi tế bào thính giác thần tốc thường để lại tác dụng phụ co giật màng nhĩ thứ cấp đấy ạ.''

Lời nhắc nhở đầy tính chuyên môn của cô em gái hóa học khiến cô chị lập tức nuốt khan, từ bỏ ngay ý định điên rồ.

Lúc này, nàng Sói đồng cỏ đã hoàn toàn phát điên vì tiếng ve sầu, vừa ôm đầu giãy giụa vừa gào thét: ``\textbf{Tại sao ở đây lại ồn ào đến mức kinh dị thế này cơ chứ!?}''

Tôi thở dài thườn thượt, giơ ngón tay chỉ thẳng về phía chiếc tủ đồ cá nhân đặt ở góc phòng: ``Đoán xem nguyên nhân từ đâu ra.''

Ngay khi ngón tay tôi vừa trỏ tới, cánh cửa tủ bỗng nhiên bật mở với một tiếng *CẠCH* giòn giã.

Bước ra từ bên trong là Haato, với nụ cười tươi rói trên môi, hoàn toàn không hề tỏ ra chút nao núng hay khó chịu nào trước mớ âm thanh chát chúa kia: ``Bởi vì chị đang nuôi ve sầu ngay trong văn phòng này đấy!''

Cả căn phòng chết lặng trong giây lát.

Nghe thì có vẻ phi lý đến mức nực cười, nhưng trước mắt chúng tôi lại là một sự thật không thể chối cãi. Trên tay cô nàng Trung học đang ôm khư khư một chiếc hộp nhựa trong suốt cỡ lớn. Bên trong đó\dots\ là cả một ổ ve sầu khổng lồ, lúc nhúc hàng chục con đang bám đầy cành cây khô và ra sức rung cánh kêu râm ran.

\img{13-10g}

Kaiho há hốc mồm, mắt trợn tròn nhìn đàn ``thú cưng'' quái đản: ``Trời đất ơi! Em đào đâu ra mà lắm ve sầu đến mức này vậy hả Haato!?''

Nàng tiểu thư tóc vàng hếch mũi lên đầy tự hào: ``Em đã mất nguyên một tuần ròng rã lặn lội trong rừng sâu để săn bắt chúng đấy ạ!''

Mồ hôi lạnh bắt đầu túa ra sau gáy tôi: ``Thế tại sao em lại tha nguyên một đàn côn trùng đến trụ sở làm việc làm cái quái gì\dots?''

Chưa kể là cái hộp nhựa của em ấy trông cũng chẳng hề kín kẽ cho cam\dots

Không thể chịu đựng thêm được nữa, Koyori lập tức rút từ trong áo khoác blouse ra một bình xịt lớn có nắp màu xanh nước biển, trên thân không dán bất kỳ nhãn mác nào.

\img{13-10f}

Trông thấy hành động đáng ngờ ấy, Kaiho giật mình: ``Cái gì nữa đây!?''

Pekora lùi lại một bước, run rẩy: ``L-Lại là thuốc độc hại não gì nữa hả, peko!?''

Nàng Sói đồng cỏ vội xua tay lia lịa, lại còn liếc mắt nhìn trộm Suzuno với vẻ sợ sệt: ``Thuốc thang gì đâu ạ! Đây chỉ là bình xịt côn trùng thông thường mua ở siêu thị thôi mà! Em thề là hàng chuẩn, không có chất cấm đâu, Suzuno-chan!''

Tôi, Kaiho và Pekora đồng loạt thở phào một tiếng: ``\dots Ồ.''

Koyori bĩu môi hờn dỗi: ``Mọi người làm cái vẻ mặt đó là sao chứ!? Bộ trong mắt mọi người em chỉ toàn mang mấy thứ vũ khí hủy diệt thôi hay sao!?''

Tôi xoa xoa cằm: ``Tại cái tiền án tiền sự của em nó dày quá rồi, không ai dám tin là em mang đồ của người bình thường cả.''

Dẫu vậy, nhìn chiếc bình xịt trên tay Koyori, tôi liền xua tay ngăn lại: ``Cơ mà\dots\ dùng bình xịt hóa chất trong phòng kín thế này thì dã man quá. Chưa kể khói độc lan ra cả lũ lại phải đi cấp cứu thì khổ.''

Fubuki cũng yếu ớt lên tiếng: ``Đúng thế đấy\dots\ Sao chúng ta không đem thả chúng ra ngoài thiên nhiên đi? Trông mấy chú ve sầu trong hộp cũng có vẻ mệt mỏi và kiệt sức lắm rồi kìa!''

Kaiho cũng gật đầu đồng tình: ``Nhốt cả đàn ve đông đúc vào một cái hộp chật chội như thế, bảo sao chúng không gào thét đòi tự do cho bằng được.''

Quả đúng như lời họ nói, bên trong hộp nhựa, nhiều con ve sầu đã bắt đầu lảo đảo, đập cánh một cách yếu ớt vì thiếu dưỡng khí.

Haato ôm chiếc hộp, gương mặt thoáng nét đắn đo, lưỡng lự. Nhìn ba cô nàng tai thú đang quằn quại chống chọi với làn sóng âm trong cơn tuyệt vọng, cuối cùng cô nàng cũng mủi lòng.

Thở dài một hơi đầy tiếc nuối, Haato cúi đầu chấp thuận: ``Thôi được rồi\dots\ Em sẽ thả chúng đi vậy.''

Nói rồi, cô nàng đưa tay chạm vào chốt nắp hộp.

Linh tính mách bảo điều chẳng lành, tôi vội vã hét toáng lên: ``\textbf{Khoan đã! Đừng có mở nắp ở ngay trong này—!}''

Nhưng đã quá muộn.

*TÁCH!*

Chiếc nắp hộp nhựa bật tung ra.

Hàng chục, thậm chí hàng trăm con ve sầu được giải phóng lập tức túa ra như ong vỡ tổ. Chúng bay tán loạn, đập cánh phần phật và lao vùn vụt khắp mọi ngóc ngách của căn phòng, kéo theo những tiếng râm ran chát chúa vang dội gấp mười lần ban nãy.

``\textbf{Trời đất ơi! Sao lại thả ngay tại chỗ này hả trời!?}'' tôi ôm đầu thét lên.

Haato chớp mắt ngây thơ: ``Ơ kìa, thì anh chị bảo em thả chúng nó ra còn gì?''

Fubuki gào lên trong nước mắt: ``\textbf{Ý là đem ra ngoài cửa sổ mà thả chứ ai bảo thả trong phòng!!}''

Kaiho cuống cuồng vơ lấy một tập hồ sơ: ``Mau tìm vợt hay bất cứ thứ gì gom chúng lại đi!!''

Cả văn phòng lập tức rơi vào một cảnh hỗn loạn chưa từng có.

Tôi và Kaiho thì vốn dĩ chẳng sợ gì ba cái thứ sâu bọ này, mặc cho lũ ve sầu lao tới bâu kín cả người. Thế nhưng, điều trớ trêu là dường như đàn ve sầu lại có một sự ``ưu ái'' kỳ lạ dành riêng cho cô bạn Tổng quản của tôi.

Hàng loạt con ve nâu to béo bay vù vù tới, đua nhau đáp thẳng lên bờ ngực đẫy đà của Kaiho, thậm chí có vài con bạo dạn còn chui tọt vào trong cổ áo và đậu xuống những chỗ hiểm hóc.

``Đông quá thể đáng!'' tôi vừa gạt một con ve trên vai vừa kêu lên.

``M-Mồ!! Sao các ngươi cứ nhắm vào ngực ta mà bâu thế hả-- Á á á! Chỗ đó không được chạm vào!!'' Kaiho đỏ bừng mặt, vừa nhảy dựng lên vừa dùng hai tay đập loạn xạ để xua đuổi.

Tôi nhịn không được, buột miệng trêu chọc: ``\dots Xem ra lũ ve này cũng biết chọn chỗ đắc địa để làm tổ đấy nhỉ.''

Kaiho thẹn quá hóa giận, hét toáng lên: ``\textbf{Ngậm miệng lại ngay cho em và mau giúp bắt chúng nó đi!! Kyaah!! Mau cút ra khỏi áo ta!!}''

Thế nhưng, so với nỗi khổ của em ấy, tình cảnh của những người còn lại mới thực sự là địa ngục trần gian.

Một con ve sầu to tướng bay thẳng tới và đáp cái bẹp ngay trên chóp mũi của Pekora, khiến cô nàng thỏ thét lên một tiếng kinh hoàng rồi ngã ngửa ra sau. Cùng lúc đó, hai con ve khác đáp xuống hai bên tai cáo của Fubuki, vô tư cất giọng rung cánh inh ỏi khiến nàng Cáo trắng run rẩy ôm đầu khóc không ra nước mắt.

\begin{figure}[H]
  \centering
  \begin{subfigure}[t]{0.45\textwidth}
    \centering
    \includegraphics[width=8cm]{../../../../Image/Book13/13-10h1.png}
  \end{subfigure}
  \quad
  \begin{subfigure}[t]{0.45\textwidth}
    \centering
    \includegraphics[width=8cm]{../../../../Image/Book13/13-10h2.png}
  \end{subfigure}
\end{figure}

Về phần Koyori, tình cảnh của em ấy còn bi hài hơn gấp bội.

Quá hoảng loạn và bế tắc, nàng Sói hồng quyết định trèo tót lên chậu cây cảnh phong thủy ở góc phòng. Hai tay hai chân em ấy ôm chặt lấy thân cây, rồi bắt đầu\dots\ giả vờ làm ve sầu để hòa nhập với đồng loại.

\img{13-10i}

``Oe oe, oe oe, oe oe\dots!''

Cái chất giọng nhái tiếng ve sầu của em ấy nghe vừa the thé, vừa thảm thiết chẳng khác nào tiếng em bé đang khóc đòi sữa. Nói ra thế này thì giống như đang cười trên sự đau khổ của người khác, nhưng đây thực sự là lần đầu tiên tôi được chứng kiến một nhà khoa học lại đi làm trò ngớ ngẩn đến mức này. Hoặc giả em ấy ngây thơ nghĩ rằng làm vậy thì lũ ve sẽ tưởng em là đồng loại mà tha cho chăng? Nếu đúng là thế thì cái biện pháp này hoàn toàn vô dụng.

Tôi đứng hình nhìn em ấy, cạn lời: ``Em kêu còn to hơn cả lũ ve thật nữa đấy Koyo-chan\dots\ Nghe như tiếng trẻ con khóc vậy. Định gia nhập bộ tộc ve sầu luôn hay sao?''

Kaiho cũng dở khóc dở cười quát với lên: ``\textbf{Trèo xuống mau em ơi! Xuống đây mà giúp chị bắt chúng nó chứ!}''

``Đông quá! Này Haato-chan, mau làm gì đó đi chứ!'' Tôi giục.

Fubuki gào lên: ``Tại sao bà lại thả bọn chúng trong nhà cơ chứ hả Haato!?''

Pekora mếu máo: ``Bắt lấy chúng đi mà, peko!''

Tôi và Kaiho đồng thanh đáp trả: ``Bọn anh/chị đang cố đây này!!''

Nàng Cáo trắng và nàng Thỏ hoảng loạn chạy tán loạn khắp phòng hòng thoát khỏi sự truy đuổi của đàn côn trùng. Trong khi đó, nàng tiểu thư tóc vàng vẫn ung dung cầm tách trà lúa mạch mát lạnh, điềm nhiên bước lại gần và cất lời giữa cơn hỗn loạn: ``Mọi người biết không? Ve sầu sau khi lột xác chỉ có thể sống được vỏn vẹn trong vòng một tuần thôi. Tuổi đời của chúng ngắn ngủi thật đấy\dots''

\img{13-10j}

Câu nói bâng quơ ấy của Haato vang lên giữa tiếng ve râm ran.

Thế nhưng, ngay khoảnh khắc những từ ngữ ấy lọt vào tai, bước chân của Suzuno bỗng nhiên khựng lại.

\dots ``\textit{Tuổi đời của chúng ngắn ngủi thật đấy.}''

Tôi chợt nhận ra bóng lưng nhỏ nhắn của em gái mình khẽ run rẩy.

Đôi mắt dị sắc màu đỏ và lam ngọc của Suzuno chùng xuống, phủ một màn sương mờ mịt. Nỗi đau của sự chia ly, của những sinh mệnh mong manh ngắn ngủi bất chợt bị khơi dậy dữ dội trong lòng cô bé.

Chỉ mới ba tuần trôi qua kể từ cái ngày định mệnh ấy. Cái ngày mà căn nhà của chúng tôi vĩnh viễn mất đi tiếng sủa rộn rã, mất đi bóng dáng của ba mẹ con nhà Chó. Mùa hè năm ngoái, Miyako, Toa và Kiri vẫn còn tíu tít chạy nhảy quanh sân, háo hức ngước nhìn những chú ve sầu trên cành cây với ánh mắt tò mò thơ ngây. Giờ đây, tiếng ve mùa hạ vẫn cất lên râm ran khắp chốn, nhưng những người bạn nhỏ bốn chân thân thương ấy\dots\ đã chẳng còn ở bên cạnh chúng tôi nữa.

Một giọt nước mắt trong suốt khẽ lăn dài trên gò má trắng ngần của Suzuno.

Trái tim tôi như thắt lại.

Quên đi lũ ve sầu đang bay loạn xạ, tôi lập tức bước nhanh tới bên cạnh Suzuno, quỳ một chân xuống và nhẹ nhàng ôm lấy bờ vai bé nhỏ của em gái. Kaiho cũng nhận ra sự biến chuyển trong tâm trạng của em gái. Cô vội vã tiến lại gần, đặt bàn tay ấm áp lên đỉnh đầu cô bé, dịu dàng xoa nhẹ: ``Suzu\dots\ Không sao đâu em.''

Tôi đưa tay khẽ gạt đi giọt nước mắt đọng nơi khóe mi của em, cất giọng trầm ấm, chan chứa yêu thương: ``Đừng buồn nữa nhé, Suzuno. Dù sinh mệnh có ngắn ngủi đến đâu, những kỷ niệm đẹp đẽ mà chúng ta đã cùng nhau trải qua sẽ mãi mãi không bao giờ biến mất đâu em.''

Suzuno cắn nhẹ môi, ngước nhìn tôi và Kaiho. Hơi ấm từ bàn tay của người anh trai và sự quan tâm chân thành của người chị đồng nghiệp dường như đã xoa dịu đi phần nào vết thương sâu thẳm trong lòng cô bé. Suzuno hít một hơi thật sâu, đưa tay quệt ngang mi mắt rồi khẽ gật đầu: ``V-Vâng ạ\dots\ Em xin lỗi vì đã làm Onii-sama và Onee-sama phải bận lòng. Em không sao rồi ạ.''

Cùng lúc đó, tiếng kêu cứu thất thanh của Pekora và Fubuki lại kéo chúng tôi về với thực tại: ``\textbf{Eek! Có con ve sầu đậu lên lọn tóc em này, peko!}'' / ``\textbf{Con ve sầu nâu đang đậu lên sát mắt em nữa!!}''

\begin{figure}[H]
  \centering
  \begin{subfigure}[t]{0.45\textwidth}
    \centering
    \includegraphics[width=8cm]{../../../../Image/Book13/13-10k1.png}
  \end{subfigure}
  \quad
  \begin{subfigure}[t]{0.45\textwidth}
    \centering
    \includegraphics[width=8cm]{../../../../Image/Book13/13-10k2.png}
  \end{subfigure}
\end{figure}

Koyori vẫn bám trên cây, kêu inh ỏi: ``\textbf{Oe oe, oe oe, oe oe\dots!}''

Fubuki và Kaiho cùng la lên: ``Chạy đi! Chạy đi mau!!''

Tôi quát to với Haato: ``Đừng có đứng đó uống trà nữa Haato-chan! Mau khẩn trương giúp bọn này bắt lại chúng nó đi!''

Nhận thấy tình hình trong phòng khách đã hoàn toàn vỡ trận, tôi lập tức chỉ tay về phía cánh cửa sắt cách âm ở cuối hành lang: ``Fubuki! Pekora! Koyori! Cả ba em mau chạy vào trốn tạm trong phòng thu âm đi!! Mau lên!!''

Như vớ được chiếc cọc cứu sinh giữa dòng nước lũ, Fubuki, Pekora và Koyori (vừa tuột khỏi thân cây) chẳng kịp suy nghĩ gì thêm, cắm đầu cắm cổ lao thục mạng về phía phòng thu âm. Cả ba lao vọt vào trong rồi dùng hết sức bình sinh đóng sập cánh cửa sắt lại.

\textit{RẦM!}

Tiếng đóng cửa vang lên khô khốc. Giờ đây, chỉ còn lại tôi, Kaiho, Haato và Suzuno ở lại tiền tuyến để giải quyết hậu quả cho bầy ve sầu ngỗ nghịch.

\pov{-}

\begin{center}
  \uline{Bên trong phòng thu âm.}
\end{center}

Nàng Thỏ và nàng Cáo trắng ngồi bệt xuống sàn, thở hổn hển như vừa hoàn thành xong một chặng đua marathon sinh tử. Cả hai ôm ngực, mồ hôi vã ra như tắm nhưng nét mặt lộ rõ vẻ nhẹ nhõm vì đã tạm thời thoát khỏi cơn ác mộng mang tên ve sầu. Trong khi đó, Koyori đang tò mò đưa mắt dáo dác nhìn quanh phòng.

\img{13-10l}

Đây là một căn phòng cách âm chuyên dụng cực kỳ kín đáo. Toàn bộ bốn bức tường đều được ốp xốp tiêu âm dày đặc, ở giữa chỉ kê duy nhất một chiếc bàn lớn cùng một chiếc ghế xoay, góc phòng là chiếc tủ đựng thiết bị trống rỗng.

Thế nhưng, có một vấn đề chết người mà cả ba người họ đã hoàn toàn không lường trước được: căn phòng này hoàn toàn không có lấy một chiếc máy điều hòa nhiệt độ, thậm chí đến một chiếc quạt máy phe phẩy cũng không có nốt. Về cơ bản, giữa cái nóng như thiêu như đốt của mùa hè tháng Tám, căn phòng kín mít không một kẽ hở này chẳng khác nào một chiếc lò bát quái khổng lồ.

Pekora vuốt mồ hôi trên trán, run rẩy hỏi: ``Chúng ta nên làm gì bây giờ đây, peko\dots?''

Fubuki tựa lưng vào tường, uể oải: ``Ta có thể ra ngoài được chưa nhỉ\dots?''

Koyori tò mò nghiêng đầu: ``Ở ngoài kia tình hình thế nào rồi ạ?''

Nàng Thỏ tóc xanh lững thững bước tới gần cánh cửa sắt, cẩn thận hé mở một khe nhỏ rộng chừng một ngón tay để thám thính tình hình.

Ngay tức khắc, một luồng âm thanh kinh hoàng gồm tiếng ve sầu kêu râm ran đinh tai nhức óc, hòa lẫn với những tiếng gầm rú và tiếng la hét bắt ve của cặp đôi Tổng quản lý ngoài kia ập thẳng vào màng nhĩ cô.

Nàng Thỏ hoảng hồn đóng sầm cửa lại, hai tay ôm chặt lấy ngực, đầu lắc nguầy nguậy như một chiếc máy giặt bị hỏng, mặt mày tái mét không còn một giọt máu: ``Không được, không được, tuyệt đối không được đâu, peko!! Ngoài kia vẫn là địa ngục trần gian đấy!!''

Nàng Cáo trắng sững sờ, hai mắt mở to trong tuyệt vọng: ``Vậy là chúng ta sẽ phải tiếp tục mắc kẹt ở trong cái phòng này ư\dots?''

Pekora gật đầu lia lịa: ``Phải chờ cho đến khi Kuon-senpai và Kaiho-senpai gom hết lũ ve đó lại thôi, peko\dots''

Đúng lúc bầu không khí đang căng thẳng tột độ, nàng Sói đồng cỏ bỗng nhiên sà tới. Koyori vênh mặt lên, nhoẻn miệng cười một cách đầy đắc thắng rồi cất giọng trêu ngươi đàn chị: ``À, thế làm sao mà à\~{}?''

*BỐP!*

Một cú tát trời giáng không một chút do dự từ Fubuki giáng thẳng vào má Koyori, kèm theo tiếng gầm đầy phẫn nộ của nàng Cáo trắng: ``\textbf{Em nên thấy may mắn vì chúng ta còn sống sót đi con Sói kia!!}''

Cô nàng Nhà khoa học ôm một bên má đỏ ửng hằn rõ dấu tay, rơm rớm nước mắt rít lên: ``C-Cái câu đó được xếp hạng tám trong danh sách `Những câu nói mà đời người nhất định phải nói một lần' đấy ạ\dots!''

Pekora tặc lưỡi, khoanh tay ngao ngán: ``Dám chắc là nếu Kuon-senpai mà nghe thấy câu đó, anh ấy sẽ gõ đầu em và bảo: `Bảng xếp hạng tào lao vô bổ gì thế hả?!', peko.''

\txtbx[1]{
  \textbf{Ghi chú thêm về bảng xếp hạng kỳ quặc của Koyori:}
  \begin{itemize}
    \item \textbf{Hạng 4:} ``Giá như chúng ta gặp nhau sớm hơn, có lẽ ta đã có thể trở thành bạn\dots''
    \item \textbf{Hạng 5:} ``Xin lỗi! Cho hỏi năm nay là năm nào vậy!?''
    \item \textbf{Hạng 6:} ``Cứ để đó cho mình! Mọi người mau chạy đi!!''
    \item \textbf{Hạng 7:} ``Tôi sẽ bao trọn gói toàn bộ chỗ này.''
    \item \textbf{Hạng 9:} ``Cuộc vui nào rồi cũng sẽ phải tàn thôi\dots''
    \item \textbf{Hạng 10:} (Dùng khi đang bẻ khóa hệ thống) ``Nào nào, bé ngoan nghe lời đi nào\~{}''
  \end{itemize}
}

\img{13-10m}

Trái ngược với vẻ bực dọc ban đầu, Fubuki - vốn là một người có máu hài hước và rất thích trêu đùa - bỗng nhiên đổi giọng, hớn hở xúi giục: ``Ồ, vậy sao? Nghe thú vị đấy chứ! Hay là em đăng bài lên Twitter luôn đi?''

Nghe tiền bối gợi ý, nàng Sói đồng cỏ lập tức quên béng đi cơn đau trên má. Em ấy hào hứng rút chiếc điện thoại thông minh ra, thoăn thoắt gõ ngay một dòng tweet với đúng câu nói đùa vừa rồi, không quên đính kèm thêm một chữ ``草'' (cười ẻ / haha) to đùng ở phía sau.

Fubuki nhìn theo, bật cười khúc khích: ``Em thực sự thích thú đến thế cơ đấy.''

Pekora thở dài ngao ngán: ``Nhưng mà em làm trò đó sai thời điểm rồi, peko\dots''

\begin{center}
  \uline{15 phút sau.}
\end{center}

Giữa một không gian kín bưng, không có lấy một luồng gió thoảng, lượng nhiệt tích tụ từ thân nhiệt và ánh nắng trưa hè bắt đầu biến căn phòng thành một chiếc lò nung thực thụ.

Nhiệt độ trong phòng tăng vọt lên từng phút một.

Mồ hôi bắt đầu đổ ra như suối, ướt đẫm mái tóc trắng và bộ trang phục của Fubuki. Nàng Cáo thở dốc, giọng run rẩy than vãn: ``Chết rồi\dots\ Chúng ta phải làm gì bây giờ đây\dots? Ta đã ở trong cái lò này quá lâu rồi\dots''

Cô gắng gượng quay đầu sang nhìn Pekora.

Đập vào mắt cô là cảnh tượng nàng Thỏ dường như đã bị ``chiên chín não bộ'' do quá nhiệt. Pekora mắt đảo tròn lờ đờ, hai chân loạng choạng vừa nhảy múa những động tác kỳ quặc vừa cất giọng hát nghêu ngao trong cơn mê sảng: ``Dàn đồng ca mùa hạ\~{} Ve trong phòng suốt ngày\~{} Mặt sàn tràn tiếng nhạc, dậy nghe nào peko\~{} Pe pe pe, peko\~{} Peko pe, peko\~{} Pe pe pe, peko\~{} Peko pe, peko\~{}!!''

\img{13-10n}

Fubuki kinh hãi quay đầu sang hướng ngược lại để cầu cứu Koyori.

Thế nhưng, nàng Sói đồng cỏ lúc này thậm chí còn không còn giữ nổi hình dáng con người. Cả cơ thể em ấy đã hoàn toàn chảy xệ ra sàn nhà như một bãi thạch màu hồng nhão nhoét, miệng liên tục lẩm bẩm trong vô thức: ``Mayonnaise\dots\ mayonnaise tan chảy rồi\dots''

\img{13-10o}

Nàng Cáo trắng chỉ còn biết nuốt một ngụm nước bọt khan đắng, bất lực thốt lên: ``Koyorin thì\dots\ cũng bị đo ván hoàn toàn rồi\dots''

Đến lúc này, bản thân cô cũng không còn đủ sức lực để chống chọi lại mức nhiệt độ khắc nghiệt thêm một giây phút nào nữa. Hai mắt cô mờ dần, đầu óc quay cuồng dữ dội. Cô quỵ ngã xuống sàn, toàn thân buông lơi và chìm dần vào bóng tối vô tận: ``Mình cũng\dots\ không thể\dots\ chịu nổi nữa rồi\dots''

Cứ như vậy, cả ba cô gái đã tự nhốt mình trong căn phòng ngột ngạt này suốt một khoảng thời gian dài. Một người thì hoá tâm thần nhảy múa mê sảng, một người thì bị nhiệt độ nung chảy thành vũng thạch, còn một người thì ngất lịm đi vì kiệt sức.

Tất cả thảm cảnh ấy chỉ bắt nguồn từ những tiếng ve sầu của mùa hè. Họ không bị tra tấn bởi thứ âm thanh inh ỏi bên ngoài thì lại bị sức nóng khủng khiếp của giữa hè đánh gục hoàn toàn\dots

%==========================================TRUYỆN PHỤ=========================================
\omk

\pov{Hikawa Kuon}

``Lời ve ngân da diết\dots''

``Se sợi chỉ âm thanh\dots''

``Khâu những đường rạo rực vào nền\dots\ mây biếc xanh\dots''

Cả ba cô gái nằm la liệt trên mặt sàn gỗ của văn phòng, người này gác chân lên người kia, miệng vẫn còn thều thào lẩm bẩm những câu thơ về ve sầu trong cơn mê man chưa dứt.

Trước mặt họ, Haato đang cầm chiếc quạt nan trên tay, liên tục phẩy mạnh để tạo ra những luồng gió mát rượi. Suzuno thì nhanh nhẹn mang tới một chậu nước ướp đá lạnh cùng những chiếc khăn bông trắng tinh, cẩn thận vắt khô rồi nhẹ nhàng đắp lên trán và sau gáy của từng tiền bối. Bên cạnh cô là ba ly nước điện giải bù khoáng tự pha chế theo đúng tỉ lệ y học an toàn, tỏa ra làn hơi mát lạnh dễ chịu.

\img{13-10p}

Tôi và Kaiho ngồi bệt bên cạnh, nhìn ba nạn nhân của đợt sốc nhiệt vừa được giải cứu mà không khỏi dở khóc dở cười.

Haato phe phẩy chiếc quạt, cảm thán: ``Hôm nay đúng là một ngày nóng kinh hoàng thật đấy nhỉ?''

Tôi dựa lưng vào tường, thở dài: ``Ừ\dots\ Lại còn đen đủi trúng ngay vào đúng lịch cắt điện bảo trì của khu phố nữa chứ. Điều hòa văn phòng cũng chẳng bật được.''

Kaiho lấy khăn lau mồ hôi trên trán, càu nhàu: ``Nếu không phải vì đàn ve sầu của em quậy phá tưng bừng thì bọn chị cũng đã dọn đồ về sớm để tránh nóng từ lâu rồi.''

Tôi quay sang lườm cô nàng tóc vàng: ``Cũng tại em cả đấy Haato. Tự dưng lại nổi hứng tha cả một bầy ve sầu từ rừng về thả rông trong trụ sở làm cái trò gì không biết.''

Nàng Trung học chắp hai tay trước ngực, chu môi làm nũng: ``Thì em đã cúi đầu xin lỗi mọi người rồi kia mà\~{}!''

Suzuno mỉm cười dịu dàng, khẽ đặt thêm một chiếc khăn lạnh lên gáy Fubuki: ``Dù sao thì Onii-sama và Onee-sama cũng đã cùng Haato-san gom hết đàn ve lại rồi thả về công viên an toàn rồi ạ. Không có con ve nào bị thương cả.''

Nghe thấy tiếng trò chuyện sột soạt xung quanh, Fubuki khẽ chớp chớp hàng mi, từ từ mở mắt ra. Thấy nàng Cáo bắt đầu cựa quậy, tôi cất tiếng hỏi: ``À, Fubu-chan, em tỉnh lại rồi đấy à?''

Nàng Cáo trắng lồm cồm bò dậy, đưa tay xoa xoa thái dương vẫn còn hơi choáng váng. Nhìn thấy khung cảnh văn phòng quen thuộc, cùng gương mặt tươi rói của Haato và nét mặt thở phào nhẹ nhõm của cặp đôi Tổng quản lý, cô nàng mới hoàn hồn nhận ra mình đã được cứu thoát khỏi chiếc lò thiêu kia.

Haato cười tươi: ``Bà cảm thấy trong người ổn hơn chưa?''

Fubuki thở hắt ra một hơi, gật đầu: ``Ừm, bà ổn rồi\dots\ Nhưng mà\dots\ hai người kia đâu rồi?''

Kaiho đưa ngón tay chỉ sang bên cạnh: ``Vẫn đang nằm im lìm như hai khúc gỗ bên cạnh em đấy thôi.''

Tôi cười khổ: ``Nhất là Koyo-chan, lúc bọn anh mở cửa ra thì trông em ấy chẳng khác nào một vũng mayonnaise bị đun sôi cả.''

Fubuki nhìn sang bên cạnh. Pekora và Koyori lúc này cũng đã dần dần lấy lại được ý thức. Nàng Thỏ xoa xoa hai tai, còn nàng Sói hồng thì co duỗi tay chân để định hình lại cơ thể.

Pekora yếu ớt thốt lên: ``Có vẻ như\dots\ peko\dots''

Koyori chớp chớp mắt, mừng rỡ: ``Bọn em đã quay trở lại làm người bình thường rồi\dots! Cảm ơn khăn lạnh của Suzuno-chan nhé!''

Suzuno mỉm cười bưng khay nước tới: ``Các chị mau uống chút nước điện giải này đi ạ. Nó sẽ giúp bù đắp lượng muối khoáng đã mất do đổ quá nhiều mồ hôi đấy ạ.''

Fubuki thấy vậy liền thở phào một hơi nhẹ nhõm, đón lấy ly nước uống một hơi cạn sạch. Nhưng rồi như sực nhớ ra điều gì, cô nàng tròn mắt đưa mắt nhìn quanh phòng khách hoàn toàn tĩnh lặng và sạch bóng côn trùng, kinh ngạc thốt lên: ``Ủa mà khoan đã\dots\ Lũ ve sầu đâu hết cả rồi!? Em nhớ rõ ràng lúc ba đứa em tháo chạy thục mạng vào phòng thu âm, ngoài này có đến cả trăm con ve bay rợp trời cơ mà!? Chẳng lẽ mọi người đã dùng bình xịt độc hại não của Koyorin tiêu diệt sạch bách rồi sao\dots?''

Pekora cũng lồm cồm nhỏm dậy, hai tai thỏ giật giật thảng thốt: ``Không lẽ Kuon-senpai và Kaiho-senpai dùng vũ lực đập bẹp dí hết chúng nó rồi, peko!?''

Tôi cười trừ, lắc đầu xua tay: ``Ai lại đi làm cái trò dã man tàn bạo ấy bao giờ. Bọn anh giải quyết chúng nó một cách hoàn toàn khoa học và hòa bình đấy chứ.''

Haato hí hửng giơ cành cây khô cùng chiếc hộp nhựa trong suốt lên khoe chiến tích: ``Tất cả là nhờ công lớn của Haachama này am hiểu tâm sinh lý của bầy ve đấy nha! Chị đã mất cả tuần lễ ăn ngủ trong rừng sâu cùng chúng nên thừa biết lũ ve sầu này cực kỳ nhạy cảm với ánh sáng tự nhiên và tiếng gõ nhịp cộng hưởng của vỏ cây!''

Suzuno mỉm cười dịu dàng, tiếp lời giải thích cặn kẽ hơn: ``Lúc đó, em đã tận dụng một chút tinh dầu khuynh diệp và tinh chất nhựa cây ngọt nhẹ có sẵn trong túi vật tư thí nghiệm mang theo, pha chế cấp tốc một dung dịch xông hơi sinh học dịu nhẹ. Làn sương tinh dầu này làm cay mắt và kích thích khứu giác của ve sầu, khiến chúng không thể bám trụ ở các góc tối, kẽ tủ hay rèm cửa, buộc phải túa ra ngoài tìm nơi thông thoáng.''

Kaiho tự hào vỗ ngực, chỉ tay về phía dãy cửa sổ sát ban công lớn: ``Kế đó, trí tuệ nhân tạo Game của Aniki quét cảm biến tầm nhiệt định vị chính xác vị trí của từng con ve một. Haato dùng cành cây mồi gõ nhịp dẫn dụ, còn anh chị mở toang toàn bộ các cánh cửa sổ hướng ra rặng cây râm mát của công viên đối diện, rồi dùng hai tấm bảng kẹp hồ sơ khổ lớn quạt tạo luồng gió mạnh, lùa toàn bộ đàn ve bay thẳng một mạch ra ngoài thiên nhiên!''

Tôi gật đầu chốt lại: ``Chỉ mất vỏn vẹn đúng mười lăm phút tác chiến. Vừa nhanh gọn, sạch sẽ, không tốn lấy một giọt hóa chất độc hại, lại không làm tổn thương hay rụng mất một chiếc cánh ve nào. Giờ thì lũ ve sầu tha hồ mà cất giọng đồng ca ngoài công viên, trả lại sự yên tĩnh cho văn phòng.''

Fubuki và Pekora nghe xong mà hai mắt tròn xoe kinh ngạc, không giấu nổi sự bái phục: ``Đúng là sự kết hợp vô tiền khoáng hậu giữa chuyên gia bắt ve tự phong và bộ ba anh em nhà Hikawa\dots\ Đỉnh thực sự đấy, peko!''

Nói đoạn, cả ba cô nàng lại khoan khoái ngả lưng nằm bệt xuống sàn, tận hưởng từng luồng gió mát rượi từ chiếc quạt nan của Haato. Trên môi của họ đều nở những nụ cười mãn nguyện như vừa từ cõi chết trở về.

Nàng Cáo trắng nhìn lên trần nhà, lẩm bẩm trong vô thức: ``Thật là một mùa hè oi bức đến ngột ngạt nhỉ\dots?''

Cả Pekora và Koyori đều đồng thanh tán thành ngay tức khắc: ``Quả đúng là như vậy đấy ạ!!''

Kaiho đứng nhìn cảnh tượng ba cô nàng nằm dài hưởng thụ, khẽ bật cười khúc khích. Rồi như sực nhớ ra điều gì, cô quay sang nhìn tôi: ``Nè, Aniki?''

Tôi quay đầu lại: ``Hửm? Có chuyện gì thế?''

Kaiho đưa tay vuốt lọn tóc mai ướt đẫm mồ hôi của mình, rồi nhìn tôi với ánh mắt đầy quan tâm: ``Anh có muốn em quạt mát cho không? Nãy giờ vật lộn bắt ve với tôi, thấy áo anh ướt sũng hết cả lưng rồi kìa.''

Tôi khẽ cười trừ: ``Thì em cũng mồ hôi mồ kê nhễ nhại kém cạnh gì anh đâu. Hôm nay đúng là một ngày siêu nóng kỷ lục mà.''

Nghe thấy lời đề nghị của em gái song sinh, Haato đang cầm quạt liền nhanh nhảu chen vào: ``Ái chà chà! Thế thì để em dùng cả hai chiếc quạt nan này phục vụ anh chị Tổng quản lý tận tình luôn nha!''

Tôi vội xua tay: ``À thôi không cần đâu, anh có thể tự lấy quạt làm được--''

Chưa kịp để tôi dứt câu, Haato đã giật phắt chiếc quạt thứ hai từ tay Kaiho, đứng chắn ngay trước mặt tôi và quạt lấy quạt để bằng cả hai tay với một tốc độ kinh hoàng.

Những luồng gió lốc mát lạnh rượi liên tục thổi tới tấp vào người tôi và Kaiho, khiến những người đang nằm bẹp dưới sàn cũng được hưởng ké luồng gió sảng khoái.

Cơn gió mát bất ngờ ập đến xua tan đi toàn bộ cái nóng nực bực bội. Tôi nhắm mắt tận hưởng, buột miệng cảm thán: ``\dots Nghĩ lại thì, anh thực sự rất cần em quạt đấy Haato-chan. Mát rượi cả người luôn.''

Lời khen vô tư ấy của tôi vô tình chạm đúng vào lòng tự ái của cô bạn đồng nghiệp bên cạnh. Kaiho lập tức phồng hai má lên giận dỗi, đôi mắt giương lên lườm tôi và Haato đầy ghen tị: ``Ăn gian quá nha Haato! Chị mới là người mở lời muốn quạt cho Aniki trước cơ mà!!''

Haato vừa quạt vừa hếch cằm trêu chọc: ``Tí nữa mới đến lượt chị làm nha! Để em phục vụ mọi người cho chu đáo đã!''

Tôi nhìn hai người họ mà dở khóc dở cười: ``Cẩn thận quạt hăng quá lại trật khớp cổ tay bây giờ đấy\dots''

Kaiho bĩu môi hờn dỗi: ``Thì anh có thèm để em quạt cho đâu mà ý với chả kiến!''

Fubuki nằm dưới sàn cười khúc khích trêu chọc: ``Kuon-senpai sướng nhất quả đất rồi nha\~{} Được tận hai mỹ nhân tranh nhau quạt mát cho cơ đấy\~{}''

Koyori cũng hùa theo, vỗ vỗ xuống sàn: ``Kuon-senpai ơi, anh cũng xuống đây nằm nghỉ ngơi với bọn em luôn đi cho mát\~{}''

Suzuno đứng bên cạnh khẽ che miệng cười trước cảnh tượng náo nhiệt của mọi người.

Nhìn quanh căn phòng, tôi thở hắt ra một hơi dài, rồi quả quyết tuyên bố: ``\dots Thật sự là công ty chúng ta nhất định phải sắm thêm vài chiếc máy phát điện dự phòng công suất lớn nữa mới được.''

Tất cả mọi người có mặt trong phòng đều đồng loạt giơ tay thốt lên đầy phấn khởi: ``Hoàn toàn tán thành\~{}\~{}!!''

\img{13-10q}

Ngoài ra, tôi cũng âm thầm ghi nhớ vào đầu kế hoạch khẩn cấp: phải lập tức cho người lắp đặt thêm một chiếc máy điều hòa riêng cùng hệ thống ống thông gió cho căn phòng thu âm chết tiệt kia. Một thiếu sót tai hại mà hai năm trước khi trang bị cơ sở vật chất ban đầu cho công ty, tôi đã vô tình bỏ quên mất.

Chốt lại thì\dots\ mùa hè năm 2022 này quả thực có quá nhiều thứ quái dị và dở khóc dở cười xảy ra nhỉ?

\end{document}